\title{Group Sequence Policy Optimization}

\begin{abstract}
We present Group Sequence Policy Optimization (GSPO), a reinforcement learning algorithm designed for the robust, effective, and efficient training of large language models. In contrast to prior methods that utilize token-wise importance ratios, GSPO employs a sequence-level approach, calculating importance ratios from entire sequence likelihoods and applying clipping, reward assignment, and optimization at the sequence granularity. Empirical results show that GSPO outperforms the GRPO algorithm in both efficiency and effectiveness, significantly enhances the stability of Mixture-of-Experts (MoE) RL training, and offers opportunities to streamline RL infrastructure. GSPO's benefits have played a key role in advancing the latest Qwen3 models.
\end{abstract}

\section{Introduction}
\label{sec:intro}

Reinforcement learning (RL) has become a central framework for enhancing the scalability of language models \citep{o1,dpsk-r1,qwq32b,qwen3}. By leveraging large-scale RL techniques, language models acquire the proficiency to address advanced tasks—such as competitive mathematics and complex programming—by engaging in extended and more nuanced reasoning.

For RL to effectively leverage increased computational resources at scale, it is essential to ensure that training remains stable and reliable. Nevertheless, leading RL methods such as GRPO \citep{grpo} suffer from pronounced instability when applied to the training of extremely large language models, frequently leading to devastating and irrecoverable collapse of the model \citep{qwen3, minimax-m1}. Such instability significantly restricts attempts to advance language model performance by extending RL-based training.

This paper reveals that the instability observed in GRPO primarily arises from an inherent misuse and breakdown of importance sampling weights within its algorithmic framework. This flaw generates substantial high-variance noise during training, which becomes more pronounced as response sequences grow longer. Additionally, the situation is exacerbated by the clipping mechanism, collectively leading to eventual model failure.

In response to these fundamental challenges, we introduce \textbf{Group Sequence Policy Optimization (GSPO)}, a novel reinforcement learning algorithm tailored for large language model training. The primary advancement of GSPO is its rigorously derived importance ratio, formulated in terms of sequence likelihood \citep{click}, which is consistent with the foundational concepts of importance sampling. Furthermore, GSPO calculates normalized rewards by measuring the advantages among several responses to a single query, thus maintaining consistency between rewards assigned at the sequence level and the optimization objective.

Our experimental results highlight that GSPO outperforms GRPO substantially in terms of training stability, efficiency, and overall effectiveness. Notably, GSPO addresses inherent stability issues typically encountered during the reinforcement learning training of large Mixture-of-Experts (MoE) architectures, thereby removing the necessity for elaborate stabilization mechanisms and paving the way for more streamlined RL system design. These strengths of GSPO directly underpin the significant performance gains achieved by the recent Qwen3 models. We anticipate that GSPO will serve as a reliable and scalable foundational algorithm, supporting ongoing progress in the large-scale reinforcement learning training of language models.

\section{Preliminaries}

\paragraph{Notation}
Throughout this work, we refer to an autoregressive language model with parameters $\theta$ as a policy, denoted by $\pi_\theta$. We represent a single query by $x$, and denote the collection of all queries as $\mathcal{D}$. For any given response $y$ corresponding to a query $x$, the probability assigned to $y$ by policy $\pi_\theta$ is expressed as $\pi_\theta (y | x)=\prod_{t=1}^{|y|} \pi_\theta (y_t | x, y_{<t} )$, where $|y|$ signifies the total number of tokens in $y$. A verifier function $r$ can evaluate a query-response pair $(x, y)$, assigning a reward value $r(x, y) \in [0, 1]$.

\paragraph{Proximal Policy Optimization (PPO)} Leveraging data drawn from the preceding policy $\pi_{\theta_\text{old}}$, PPO \citep{ppo} restricts policy modifications to remain near the original policy by introducing a clipping technique. In particular, PPO optimizes the policy based on the following objective function (excluding the KL regularization term, which is not central to the present discussion):

\begin{align}
\mathcal{J}_\text{PPO}(\theta) = \mathbb{E}_{ x \sim \mathcal{D},\, y \sim \pi_{\theta_\text{old}}( \cdot | x) }
\left[ \frac{1}{|y|} \sum_{t=1}^{|y|} 
\min \left( w_{t}(\theta) \widehat{A}_{t},  \, \mathrm{clip} \left( w_{t}(\theta), 1 - {\varepsilon}, 1 + {\varepsilon}\right) \widehat{A}_{t} \right)
\right],
\end{align}

Here, the token $y_t$'s importance ratio is expressed as $ w_{t}(\theta) = \frac{ \pi_{\theta} (y_{t} | x, y_{<t}) }{ \pi_{\theta_\text{old}} (y_{t} | x,y_{<t})} $. The advantage $\widehat{A}_{t}$ assigned to $y_t$ is computed using a separate value model, while $\varepsilon$ denotes the threshold for clipping the importance ratios.

A primary obstacle in the practical implementation of PPO is its significant dependence on the value model. In practice, the value model typically matches the policy model in terms of parameter count, resulting in substantial demands on both memory and computation. Moreover, the overall performance of the algorithm is closely tied to the precision of its value estimates. Achieving a dependable value model is already a non-trivial problem; extending this reliability to scenarios involving lengthier responses and increasingly sophisticated tasks proves to be an even more formidable issue.

\paragraph{Group Relative Policy Optimization (GRPO)} GRPO \citep{grpo} eliminates the requirement for a value model by assessing the relative advantage among responses clustered together for a given prompt. More precisely, GRPO aims to maximize the following objective:

\begin{align}
\label{equ:grpo}
\mathcal{J}_\text{GRPO}(\theta) = \mathbb{E}_{ x \sim \mathcal{D},\, \{y_i\}_{i=1}^G \sim \pi_{\theta_\text{old}}( \cdot | x) }
\left[ \frac{1}{G} \sum_{i=1}^{G} \frac{1}{|y_i|} \sum_{t=1}^{|y_i|} 
\min \left( w_{i,t}(\theta) \widehat{A}_{i,t},  \, \mathrm{clip} \left( w_{i,t}(\theta), 1 - {\varepsilon}, 1 + {\varepsilon}\right) \widehat{A}_{i,t} \right)
\right],
\end{align}

where $G$ is the number of generated responses to each query $x$ (i.e., the group size), and the importance ratio $w_{i,t}(\theta)$ and advantage $\widehat{A}_{i,t}$ of token $y_{i,t}$ are:

\begin{align}
    w_{i,t}(\theta)=\frac{ \pi_{\theta} (y_{i,t} | x, y_{i,<t}) }{ \pi_{\theta_\text{old}} (y_{i,t} | x,y_{i,<t})},\quad\ 
    \widehat{A}_{i,t} = \widehat{A}_{i} = \frac{r(x, y_i) - \mathrm{mean} \left( \{ r(x, y_i) \}_{i=1}^G \right) }{ \mathrm{std} \left( \{ r(x, y_i) \}_{i=1}^G \right) },
\end{align}

respectively, where all the tokens in $y_i$ share the same advantage as $\widehat{A}_{i}$.

\section{Motivation}
\label{sec:motivation}

The expansion of model parameters, increased sparsity (such as observed in Mixture-of-Experts models), and extended response sequences collectively require the use of large rollout batch sizes to fully exploit hardware capabilities during RL. To enhance the efficiency with which samples are used, it is a common approach to subdivide these extensive batches into several mini-batches for performing gradient updates. However, this leads to an off-policy learning context, in which the collected responses $y$ stem from a previous policy $\pi_{\theta_\text{old}}$, rather than from the current policy $\pi_\theta$ under optimization. As a result, mechanisms such as clipping—used in PPO and GRPO—are essential for preventing gradient calculations from relying excessively on samples that deviate significantly from the current policy.

Although approaches such as clipping are designed to mitigate the off-policy discrepancy, we observe a deeper flaw inherent to GRPO: \textit{the objective itself lacks a well-defined formulation}. This limitation becomes especially problematic in scenarios involving the training of large-scale models on tasks with extended responses, frequently resulting in severe model degradation. The root of this issue lies in the improper utilization of importance sampling weights within the GRPO framework. Importance sampling is fundamentally intended to approximate the expectation of a function $f$ with respect to a target distribution $\pi_\text{tar}$ by appropriately adjusting the weights of samples obtained from a behavior distribution $\pi_\text{beh}$:

\begin{align}
\mathbb{E}_{z \sim \pi_\text{tar}} \left[ f(z) \right] = \mathbb{E}_{z \sim \pi_\text{beh}} \left[ \frac{ \pi_\text{tar} (z) }{ \pi_\text{beh} (z)} \, f(z) \right].
\end{align}

Importantly, effective correction for the mismatch between distributions using the importance weight $\frac{ \pi_\text{tar} (z) }{ \pi_\text{beh} (z)}$ requires averaging across a large number of samples ($N \gg 1$) drawn from the behavior distribution $\pi_\text{beh}$.

Conversely, GRPO assigns the importance weight $\frac{ \pi_{\theta} (y_{i,t} | x, y_{i,<t}) }{ \pi_{\theta_\text{old}} (y_{i,t} | x,y_{i,<t})}$ independently at each token position $t$. This importance weighting depends only on a solitary sample $y_{i,t}$ drawn from the next-token distribution $\pi_{\theta_\text{old}}( \cdot | x, y_{i,<t})$ and thus fails to correct for distribution mismatch as intended. Rather than mitigating distribution shift, this approach injects substantial variance into the gradients used for training. This elevated noise accumulates over extended sequences and is aggravated further by the use of clipping. Our experimental findings indicate that such dynamics can precipitate a catastrophic and often irreversible collapse of the model. After such collapse, further training does not recover model performance, even if training resumes from an earlier checkpoint or after extensive efforts to optimize hyperparameters (for example, by altering clipping bounds, increasing sequence length, or modifying the RL querying strategy).

The observation highlighted above reveals a key limitation in the architecture of GRPO. Specifically, the inadequacy of token-level importance weighting underscores an essential concept: \textbf{the granularity of the optimization objective must be consistent with that of the reward}. Given that rewards are assigned for whole sequences rather than individual tokens, implementing off-policy corrections at the token level may introduce issues. Consequently, this realization prompts us to abandon token-level objectives and instead consider formulating importance weights and conducting optimization at the \textit{sequence level}.

\section{Algorithm}

\subsection{GSPO: Group Sequence Policy Optimization}

Although the token-level importance weight $\frac{ \pi_{\theta} (y_{i,t} | x, y_{i,<t}) }{ \pi_{\theta_\text{old}} (y_{i,t} | x,y_{i,<t})}$ presents issues in GRPO, our analysis reveals that for language generation tasks, the \textit{sequence-level} importance weight $\frac{ \pi_{\theta} (y | x) }{ \pi_{\theta_\text{old}} (y | x)}$ possesses a well-defined theoretical interpretation. Specifically, it quantifies the extent to which a sampled response $y$ from $\pi_{\theta_\text{old}} (\cdot | x)$ differs from $\pi_{\theta} (\cdot | x)$, closely matching sequence-level rewards. Furthermore, this metric provides a useful signal for implementing the clipping procedure.

Based on this straightforward observation, we propose the \textbf{Group Sequence Policy Optimization (GSPO)} algorithm. GSPO employs the following sequence-level optimization objective:

\begin{align}
\mathcal{J}_\text{GSPO} (\theta) =
\mathbb{E}_{ x \sim \mathcal{D},\, \{y_i\}_{i=1}^G \sim \pi_{\theta_\text{old}}( \cdot | x) }
\left[ 
\frac{1}{G} \sum_{i=1}^{G}
\min \bigg( s_{i}(\theta)\widehat{A}_{i},  \; \mathrm{clip}\left( s_{i}(\theta), 1 - {\varepsilon}, 1 + {\varepsilon} \right)\widehat{A}_{i} \bigg)
\right],
\label{equ:gspo}
\end{align}

where we adopt the group-based advantage estimation:

\begin{align}
\widehat{A}_{i} = \frac{r(x, y_i) - \mathrm{mean} \left( \{ r(x, y_i) \}_{i=1}^G \right) }{ \mathrm{std} \left( \{ r(x, y_i) \}_{i=1}^G \right) },
\end{align}

and define the importance ratio $s_{i}(\theta)$ based on sequence likelihood \citep{click}:

\begin{align}
s_{i}(\theta) = \left( \frac{ \pi_{\theta} (y_i | x) }{ \pi_{\theta_\text{old}} (y_i | x)} \right)^{\frac{1}{|y_i|}}
=
\exp \left( \frac{1}{|y_i|} \sum_{t=1}^{|y_i|} \log \frac{ \pi_{\theta} (y_{i,t} | x, y_{i,<t}) }{ \pi_{\theta_\text{old}} (y_{i,t} | x,y_{i,<t})} \right).
\end{align}

Consequently, GSPO employs clipping at the response level, rather than at the token level, to filter out excessively “off-policy” samples during gradient calculation. This strategy aligns both the reward assignment and optimization process at the level of full sequences. Importantly, we utilize length normalization in $s_{i}(\theta)$ to mitigate variance and to keep $s_{i}(\theta)$ within a consistent numerical interval. Without this normalization, significant shifts in the likelihood of a small subset of tokens could cause large oscillations in the sequence-level importance ratio, thereby necessitating different clipping thresholds for responses of varying lengths. Additionally, it is important to mention that the magnitude of clipping intervals used in GSPO generally contrasts with those in earlier methods (such as GRPO), which stems from the differing formulations of the importance ratios in each case.

\subsection{Gradient Analysis}
\label{subsec:gradient}

We can derive the gradient of the GSPO objective as follows (clipping is omitted for brevity):

\begin{align}
\nabla_{\theta} \mathcal{J}_\text{GSPO} (\theta)
=&\ 
\nabla_{\theta} \mathbb{E}_{ x \sim \mathcal{D},\, \{y_i\}_{i=1}^G \sim \pi_{\theta_\text{old}}( \cdot | x) }
\left[ 
\frac{1}{G} \sum_{i=1}^{G} s_{i}(\theta) \widehat{A}_{i}
\right] \\
= &\ 
\mathbb{E}_{ x \sim \mathcal{D},\, \{y_i\}_{i=1}^G \sim \pi_{\theta_\text{old}}( \cdot | x) }
\left[ 
\frac{1}{G} \sum_{i=1}^{G}
s_{i}(\theta) \widehat{A}_{i}
\, \nabla_{\theta} \log s_{i}(\theta)
\right] \\
=&\ 
\mathbb{E}_{ x \sim \mathcal{D},\, \{y_i\}_{i=1}^G \sim \pi_{\theta_\text{old}}( \cdot | x) }
\left[ 
\frac{1}{G} \sum_{i=1}^{G} 
\left( \frac{ \pi_{\theta} (y_i | x) }{ \pi_{\theta_\text{old}} (y_i | x)} \right)^{\frac{1}{|y_i|}} \widehat{A}_{i} 
\, \frac{1}{|y_i|} \sum_{t=1}^{|y_i|} \nabla_{\theta} \log \pi_{\theta} (y_{i,t} | x, y_{i,<t}) 
\right].
\label{equ:gspo_grad}
\end{align}

For comparison, the gradient of the GRPO objective is as follows (note that $\widehat{A}_{i,t} = \widehat{A}_{i}$):

\begin{align}
\nabla_{\theta} \mathcal{J}_\text{GRPO} (\theta) 
=&\
\nabla_{\theta} \mathbb{E}_{ x \sim \mathcal{D},\, \{y_i\}_{i=1}^G \sim \pi_{\theta_\text{old}}( \cdot | x) }
\left[ \frac{1}{G} \sum_{i=1}^{G} \frac{1}{|y_i|} \sum_{t=1}^{|y_i|} 
w_{i,t}(\theta) \widehat{A}_{i,t}
\right] \\
=&\
\mathbb{E}_{ x \sim \mathcal{D},\, \{y_i\}_{i=1}^G \sim \pi_{\theta_\text{old}}( \cdot | x) }
\left[
\frac{1}{G} \sum_{i=1}^{G} \widehat{A}_{i} 
\cdot \frac{1}{|y_i|} \sum_{t=1}^{|y_i|} 
\frac{ \pi_{\theta}(y_{i,t} | x, y_{i,<t}) }{ \pi_{\theta_\text{old}}(y_{i,t} | x, y_{i,<t}) }
\nabla_{\theta} \log \pi_{\theta}(y_{i,t} | x, y_{i,<t})
\right].
\end{align}

The primary difference between GSPO and GRPO concerns \textit{the manner in which they assign weights to the gradients of token log likelihoods}. GRPO employs an ``importance weight'' for each token, given by $\frac{ \pi_{\theta} (y_{i,t} | x, y_{i,<t}) }{ \pi_{\theta_\text{old}} (y_{i,t} | x,y_{i,<t})}$, to modulate their influence. These weights, which are distributed within the ranges $(0, 1+\varepsilon]$ when $\widehat{A}_{i} >0$ or $[1-\varepsilon, +\infty)$ when $\widehat{A}_{i} < 0$, can differ significantly. Such variance in weighting is nontrivial and, over time, can accumulate, resulting in potentially erratic outcomes during model optimization. Conversely, GSPO utilizes uniform weighting across all tokens in a given response, thus removing the instability introduced by the disparate weights of GRPO.

\subsection{GSPO-token: A Token-level Objective Variant}

In situations such as multi-turn RL, it can be beneficial to modify the advantage at a more granular level than the entire sequence. For this purpose, we propose a variant of GSPO at the token level, referred to as \textbf{GSPO-token}, enabling the adjustment of advantages on a per-token basis:

\begin{align}
\mathcal{J}_\text{GSPO-token}(\theta)
=
\mathbb{E}_{ x \sim \mathcal{D},\, \{y_i\}_{i=1}^G \sim \pi_{\theta_\text{old}}( \cdot | x) }
\left[ 
\frac{1}{G} \sum_{i=1}^{G} \frac{1}{|y_i|} \sum_{t=1}^{|y_i|} 
\min \Bigg( s_{i,t}(\theta) \widehat{A}_{i,t},\  
\mathrm{clip} \left( s_{i,t}(\theta), 1 - \varepsilon, 1 + \varepsilon \right) \widehat{A}_{i,t} \Bigg) 
\right],
\label{equ:gspo-token}
\end{align}

where

\begin{align}
s_{i,t}(\theta) 
= \mathrm{sg} \left[ s_{i}(\theta) \right]  \cdot \frac{ \pi_{\theta} (y_{i,t} | x, y_{i,<t}) }{ \mathrm{sg} \left[ \pi_{\theta} (y_{i,t} | x, y_{i,<t}) \right] },
\end{align}

and $\mathrm{sg}[\cdot]$ denotes only taking the numerical value but stopping the gradient, corresponding to the \texttt{detach} operation in PyTorch. The gradient of GSPO-token can be derived as:

\begin{align}
\nabla_{\theta} \mathcal{J}_\text{GSPO-token}(\theta)
=&\ 
\nabla_{\theta} \mathbb{E}_{ x \sim \mathcal{D},\, \{y_i\}_{i=1}^G \sim \pi_{\theta_\text{old}}( \cdot | x) }
\left[ 
\frac{1}{G} \sum_{i=1}^{G}
\frac{1}{|y_i|} \sum_{t=1}^{|y_i|} 
s_{i,t}(\theta) \widehat{A}_{i,t}
\right] \\
=&\ 
\mathbb{E}_{ x \sim \mathcal{D},\, \{y_i\}_{i=1}^G \sim \pi_{\theta_\text{old}}( \cdot | x) }
\left[ 
\frac{1}{G} \sum_{i=1}^{G} s_i (\theta) \cdot \frac{1}{|y_i|} \sum_{t=1}^{|y_i|} 
\widehat{A}_{i,t} \frac{ \nabla_{\theta} \pi_{\theta} (y_{i,t} | x, y_{i,<t}) }{ \pi_{\theta} (y_{i,t} | x, y_{i,<t}) }
\right] \\
=&\ 
\mathbb{E}_{ x \sim \mathcal{D},\, \{y_i\}_{i=1}^G \sim \pi_{\theta_\text{old}}( \cdot | x) }
\left[
\frac{1}{G} \sum_{i=1}^{G} \left( \frac{ \pi_{\theta} (y_i | x) }{ \pi_{\theta_\text{old}} (y_i | x)} \right)^{\frac{1}{|y_i|}}
\cdot \frac{1}{|y_i|} \sum_{t=1}^{|y_i|}
\widehat{A}_{i,t} \nabla_{\theta} \log \pi_{\theta} (y_{i,t} | x, y_{i,<t})
\right].
\label{equ:gspo-token_grad}
\end{align}

Observe that the expression $\frac{ \pi_{\theta} (y_{i,t} | x, y_{i,<t}) }{ \mathrm{sg} \left[ \pi_{\theta} (y_{i,t} | x, y_{i,<t}) \right] }$ evaluates to 1, which implies that, in practice, $s_{i,t}(\theta)$ and $s_{i}(\theta)$ yield the same numerical result.

When we examine Equation~\eqref{equ:gspo} versus~\eqref{equ:gspo-token}, as well as Equation~\eqref{equ:gspo_grad} compared to~\eqref{equ:gspo-token_grad}, we find that GSPO-token and GSPO are equivalent in terms of optimization objectives, the conditions for clipping, and their theoretical gradients, as long as the advantage is kept constant across every token within the response $y_i$ (meaning $\widehat{A}_{i,t} = \widehat{A}_{i}$ for all $t$). However, GSPO-token extends further flexibility, since it allows the advantage values to be set individually for each token.

\section{Experiments and Discussion}

\subsection{Empirical Results}

Our experiments utilize a cold-start model, specifically fine-tuned from Qwen3-30B-A3B-Base. We present both the training reward trajectories and model performance metrics on three benchmarks: AIME'24 (evaluated using average Pass@1 over 32 samples), LiveCodeBench (202410-202502, assessed by average Pass@1 over 8 samples), and CodeForces (using Elo Rating). During reinforcement learning training, the rollout data in each batch is divided into four mini-batches to perform gradient updates. For GSPO, the left and right clipping thresholds in Equation~\eqref{equ:gspo} are chosen as 3e-4 and 4e-4, respectively. We use GRPO as the baseline method, with its corresponding clipping ranges in Equation~\eqref{equ:grpo} set to 0.2 and 0.27, values that have been meticulously selected to provide an equitable comparison. Importantly, we highlight that normal MoE RL convergence with GRPO requires the Routing Replay training approach—which will be further discussed in \S~\ref{subsec:routing_replay}—whereas \textit{GSPO eliminates this requirement}.

Figure~\ref{fig:results} demonstrates that training with GSPO maintains stable progress at all stages. Notably, \textbf{GSPO enables ongoing enhancements in performance when additional computational resources are allocated for training, the query set is periodically refreshed, and the length of generated sequences is increased}. In addition, GSPO surpasses GRPO in terms of training efficiency, achieving higher accuracy and benchmark results given equivalent training compute and the same number of utilized queries. Lastly, the application of GSPO in RL training for the latest Qwen3 models provides strong evidence for the effectiveness of GSPO in harnessing the scaling capabilities of RL for large language models.

\subsection{Curious Observation on Clipping Fractions}

One fundamental difference between GSPO and GRPO lies in their respective clipping strategies: GSPO applies clipping at the response level, whereas GRPO operates at the token level. As illustrated in Figure~\ref{fig:clipping}, this results in GSPO clipping a vastly higher proportion of tokens—by about two orders of magnitude—compared to GRPO. This substantial gap remains unaffected even when the clipping thresholds are altered. Remarkably, even though GSPO discards a much larger number of tokens (thus utilizing fewer tokens in training or for estimating gradients), it nonetheless achieves greater training efficiency than GRPO. This somewhat paradoxical outcome suggests that GRPO's token-wise gradient estimations suffer from high noise and inefficiency in leveraging samples. In contrast, GSPO's approach, which focuses on the entire sequence, yields a more stable and potent learning signal.

\subsection{Benefit of GSPO for MoE Training}
\label{subsec:routing_replay}

\paragraph{Background}
In contrast to RL training on dense models, the sparse expert activation characteristic of MoE architectures brings distinct challenges regarding training stability. Notably, when implementing the GRPO algorithm, the phenomenon of \textit{expert-activation volatility} in MoE systems may impede proper convergence of RL training. Specifically, after performing one or more gradient steps, the subset of experts engaged for producing identical outputs can vary greatly. To illustrate, in the case of the 48-layer Qwen3-30B-A3B-Base, each RL gradient iteration leads to approximately a 10\% turnover in the experts triggered for the same rollout sample, comparing the new policy $\pi_{\theta}$ to the prior policy $\pi_{\theta_\text{old}}$.
This volatility, which intensifies in deeper MoE architectures, causes token-level importance weights $w_{i,t}(\theta) = \frac{ \pi_{\theta} (y_{i,t} | x, y_{i,<t}) }{ \pi_{\theta_\text{old}} (y_{i,t} | x,y_{i,<t})}$ to exhibit significant fluctuations, thereby compromising their reliability, as elaborated in \S~\ref{sec:motivation} and~\ref{subsec:gradient}. These instabilities ultimately obstruct successful RL training convergence.

\paragraph{Our Previous Approach} In our earlier work, we addressed this issue through the use of the \textbf{Routing Replay} training method. This approach involves recording the set of experts activated by $\pi_{\theta_\text{old}}$ and subsequently using these identical routing decisions when evaluating $\pi_{\theta}$ to calculate the importance weights $w_{i,t}(\theta) = \frac{ \pi_{\theta} (y_{i,t} | x, y_{i,<t}) }{ \pi_{\theta_\text{old}} (y_{i,t} | x,y_{i,<t})}$. As a result, the two models—$\pi_{\theta}$ and $\pi_{\theta_\text{old}}$—activate the same experts for each token $y_{i,t}$, thus maintaining alignment between their underlying networks. This alignment helps stabilize token-wise importance ratios and promotes consistent optimization of the corresponding subnetwork as gradients are updated. As illustrated in Figure~\ref{fig:routing_replay}, Routing Replay proves to be a crucial mechanism that enables successful and stable convergence during GRPO training of MoE models.

\paragraph{Benefit of GSPO} While Routing Replay allows GRPO training to converge effectively for MoE models, it does so at the cost of increased memory use and communication, as it relies on storing and reusing routing patterns. This process can restrict the effective capacity of the MoE architecture. In comparison, as depicted in Figure~\ref{fig:results}, GSPO operates independently of Routing Replay, permitting direct computation of importance ratios $s_i(\theta)$ as usual and ensuring both normal convergence and stable optimization. The principal insight lies in GSPO's emphasis on the sequence-level likelihood (that is, $\pi_{\theta} (y_i | x)$) while exhibiting robustness to the likelihoods of individual tokens (specifically, $\pi_{\theta} (y_{i,t} | x, y_{i,<t})$). Given the persistent language modeling proficiency of the MoE, the sequence likelihood remains relatively steady over time. Thus, GSPO fundamentally addresses the volatility in expert activation that plagues MoE models, thereby removing the necessity for intricate solutions such as Routing Replay. This leads to a more streamlined and robust training dynamic and enables the MoE model to fully exploit its inherent capacity without imposed limitations.

\subsection{Benefit of GSPO for RL Infrastructure}

Due to the differences in numerical precision between engines used for training (such as Megatron) and those designed for inference (such as SGLang and vLLM), it is common practice to recalculate the likelihoods of sampled responses under the prior policy $\pi_{\theta_\text{old}}$ using the training engine. Nonetheless, GSPO optimizes over sequence-level likelihoods rather than token-level ones, which, intuitively, are significantly less sensitive to precision mismatches. As a result, GSPO enables direct utilization of the likelihoods produced by the inference engine for optimization, eliminating the necessity of recomputation via the training engine. This feature is particularly advantageous in contexts like partial rollout, multi-turn RL, or within frameworks that separate training and inference workflows.

\section{Conclusion}

We introduce Group Sequence Policy Optimization (GSPO), an innovative reinforcement learning method designed for the training of large language models. By adhering to the core concept of importance sampling, GSPO establishes sequence-level importance ratios, and implements sequence-wise clipping, rewarding, and optimization strategies. Relative to GRPO, GSPO offers marked advantages in terms of training stability, sample efficiency, and overall performance, and proves especially effective in large-scale reinforcement learning scenarios with MoE models, underpinning significant advancements observed in recent Qwen3 models. Positioned as a scalable and robust foundation, GSPO enables further scaling of reinforcement learning approaches, paving the way for transformative progress in artificial intelligence.

\end{document}